\section{Preliminaries, Research Questions, and Contributions}
\label{sec:preliminaries}

The method combines executable process semantics with a statistical predictor. To make the interface between them precise, this section first defines the alignment problem and then introduces the small set of machine learning concepts used by LARA. The emphasis is on what each concept permits us to claim. These definitions also establish the research questions and the evidence needed to answer them.

\subsection{Traces, Variants, and Labeled Petri Nets}

A trace records the order of activities observed for one case. Let $\mathcal A$ be an activity alphabet and let $\sigma=\langle a_1,\ldots,a_n\rangle\in\mathcal A^*$ be a finite trace. Its length is $|\sigma|=n$; the empty trace is $\epsilon$. An event log $\mathcal E$ is a multiset of traces. A \emph{variant} is a distinct activity sequence, and its frequency is the number of cases with that sequence. LARA studies control flow: timestamps, resources, and other event attributes do not enter the scorer or alignment cost.

We represent a labeled Petri net as $SN=(P,T,F,\lambda,m_0,m_f)$. Here $P$ is a finite set of places, $T$ is a disjoint finite set of transitions, and $F\subseteq(P\times T)\cup(T\times P)$ is the arc relation. We use ordinary arcs with weight one in the formal exposition. The labeling function $\lambda:T\to\mathcal A\cup\{\tau\}$ assigns a visible activity or the invisible label $\tau$. A marking $m:P\to\N$ assigns tokens to places, with $\N$ including zero. The initial and final markings are $m_0$ and $m_f$.

For $t\in T$, its preset ${}^\bullet t=\{p:(p,t)\in F\}$ contains its input places, and its postset $t^\bullet=\{p:(t,p)\in F\}$ contains its output places. The transition is enabled at $m$ when $m(p)\geq1$ for every $p\in{}^\bullet t$. Firing it produces the marking
\begin{equation}
 m'(p)=m(p)-\ind[p\in{}^\bullet t]+\ind[p\in t^\bullet],
 \qquad p\in P,
 \label{eq:firing}
\end{equation}
written $m\xrightarrow{t}m'$. In words, firing consumes a token from every input and produces a token in every output. A transition whose label is $\tau$ still changes the marking; it simply does not represent a recorded activity. A \emph{complete firing sequence} starts at $m_0$ and reaches exactly $m_f$. We assume that at least one such sequence exists. This assumption ensures that an alignment can always be constructed by using log moves for the trace and model moves for a complete model execution.

Distinct transitions may carry the same label. We therefore distinguish the identifier $t$ from its visible label $\lambda(t)$. Figure~\ref{fig:running} illustrates why: $t_1$ and $t_2$ both explain an event labeled $b$, but only $t_1$ directly leads to $c$. Matching the label alone does not identify a model execution. Invisible transitions create a related issue: an event may become executable only after unobserved routing steps.

\begin{figure*}[!t]
\centering
\begin{tikzpicture}[x=1cm,y=1cm]
\node[paneltitle] at (0,1.85) {One observed label, two possible transitions};
\node[place] (p0) at (.3,0) {}; \fill[ink] (p0) circle (.065);
\node[transition] (ta) at (1.45,0) {$a$};
\node[place] (p1) at (2.6,0) {};
\node[transition] (t1) at (3.8,.78) {$b$};
\node[transition] (t2) at (3.8,-.78) {$b$};
\node[place] (p2) at (5,.78) {};
\node[place] (p3) at (5,-.78) {};
\node[transition] (tc) at (6.2,.78) {$c$};
\node[transition] (td) at (6.2,-.78) {$d$};
\node[place] (p4) at (7.4,0) {};
\node[silent] (tt) at (8.6,0) {$\tau$};
\node[place,double] (p5) at (9.8,0) {};
\foreach \a/\b in {p0/ta,ta/p1,p1/t1,p1/t2,t1/p2,t2/p3,p2/tc,p3/td,tc/p4,td/p4,p4/tt,tt/p5}
  \draw[arc] (\a) -- (\b);
\foreach \a/\l in {p0/p_0,p1/p_1,p4/p_4,p5/p_5}
  \node[annotation,below=3mm of \a] {$\l$};
\foreach \a/\l in {ta/t_a,t1/t_1,tc/t_c,tt/t_\tau}
  \node[annotation,above=2mm of \a] {$\l$};
\node[annotation,below=2mm of t2] {$t_2$};
\node[annotation,below=2mm of td] {$t_d$};
\node[annotation,above=2mm of p2] {$p_2$};
\node[annotation,below=2mm of p3] {$p_3$};
\node[paneltitle] at (0,-2.25) {A complete zero cost alignment for $\sigma=\langle a,b,c\rangle$};
\node[annotation,anchor=east] at (1.8,-2.95) {log};
\node[annotation,anchor=east] at (1.8,-3.65) {model};
\foreach \x/\a/\t in {2.6/a/t_a,3.7/b/t_1,4.8/c/t_c,5.9/\skipmove/t_\tau}{
  \node[alcell] at (\x,-2.95) {$\a$};
  \node[alcell] at (\x,-3.65) {$\t$};}
\node[verifybox,text width=31mm] at (8.5,-3.3) {replay accepts\\$m_f$ reached\\$c(\gamma)=0$};
\node[annotation,anchor=west] at (0,-4.5) {$[p_0]\xrightarrow{t_a}[p_1]\xrightarrow{t_1}[p_2]\xrightarrow{t_c}[p_4]\xrightarrow{t_\tau}[p_5]=m_f$};
\end{tikzpicture}

\caption{A running example of duplicate label ambiguity. The token begins at $p_0$ and must finish at $p_5$. Both $t_1$ and $t_2$ are labeled $b$, but their continuations differ. For $\langle a,b,c\rangle$, the upper branch gives a zero cost alignment. The final invisible transition must still fire even though it consumes no event. Transition identifiers, activity labels, and alignment skips have separate roles.}
\label{fig:running}
\end{figure*}

The \emph{complete visible language} $\mathcal L(SN)$ contains the label sequences of complete firing sequences after removing all $\tau$ labels. Two Petri nets are language equivalent if their complete visible languages coincide. This is a behavioral relation, not a claim that their graphs or internal choices are identical. An \emph{isomorphic renaming} is stronger: it changes identifiers through bijections on places and transitions while preserving arcs, labels, and markings. We use both relations to test whether candidate quality depends on a particular representation.

Process trees provide one source of synthetic Petri nets. A visible leaf represents an activity; sequence executes children in order, exclusive choice (XOR) selects one child, and parallel composition (AND) permits their interleavings. A loop permits repetition. The training generator uses sequence, XOR, and AND; the separate stress generator also includes loops. A fixed compiler converts these operators to Petri net fragments. We call its output \emph{canonical block} only to name this particular compilation, without asserting a universal normal form.

\subsection{Alignments: Executability and Deviation Cost}
\label{sec:alignment}

An alignment is needed because a trace may omit a model activity or contain an activity that the model cannot execute at that point. Let $\skipmove$ denote a skip on one side of an alignment. The permitted moves and their costs are shown in Table~\ref{tab:moves}. A synchronous move consumes one observed event and fires a transition with that label. A log move consumes an event without changing the marking. A model move fires a transition without consuming an event. The pair $(\skipmove,\skipmove)$ is excluded because it explains nothing.

\begin{table*}[!t]
\caption{Move vocabulary and the fixed costs used in this study. An invisible model move is an actual firing, whereas a skip means that nothing happens on that side.}
\label{tab:moves}
\centering
\begin{tabularx}{\linewidth}{@{}l l Y c@{}}
\toprule
Move & Form & Meaning & Cost\\
\midrule
Synchronous & $(a,t)$, $\lambda(t)=a$ & The event agrees with a model firing. & 0\\
Log & $(a,\skipmove)$ & The event has no corresponding firing. & 1\\
Visible model & $(\skipmove,t)$, $\lambda(t)\ne\tau$ & A modeled activity is absent from the trace. & 1\\
Invisible model & $(\skipmove,t)$, $\lambda(t)=\tau$ & Unobserved routing advances the model. & 0\\
\bottomrule
\end{tabularx}
\end{table*}

For a move sequence $\gamma$, its log projection $\pi_L(\gamma)$ removes the model components and skips; its model projection $\pi_T(\gamma)$ removes the log components and skips but retains all transition identities, including invisible transitions. The alignment is \emph{legal} when
\begin{equation}
 \pi_L(\gamma)=\sigma
 \quad\text{and}\quad
 m_0\xrightarrow{\pi_T(\gamma)}m_f,
 \label{eq:legal}
\end{equation}
and every synchronous move has matching labels. Thus a legal alignment reconstructs the observed sequence exactly and describes a complete executable model path. Merely firing enabled transitions is insufficient: a sequence can stop in the wrong marking.

Let $c(\gamma)$ be the sum of move costs. The optimal alignment cost is
\begin{equation}
 \delta(\sigma,SN)=\min_{\gamma\text{ legal for }(\sigma,SN)}c(\gamma).
 \label{eq:optimum}
\end{equation}
Under these costs, $\delta$ is the minimum number of visible insertions and deletions needed to explain the trace by a complete model execution. It does not measure a monetary or business loss. A substitution normally requires one log move and one visible model move, giving cost two. Other cost policies may be appropriate in an application, but the empirical results here concern the policy in Table~\ref{tab:moves}.

An optimum need not be unique. Independent activities can be ordered differently, and invisible routes may differ without changing the visible cost. Consequently, matching the exact teacher's move sequence is a stricter diagnostic than achieving the same cost. We will report both, but use cost equality for optimality. For a candidate verified by replay $\widehat\gamma$, its \emph{gap} is
\begin{equation}
 g=c(\widehat\gamma)-\delta(\sigma,SN)\geq0.
 \label{eq:gap}
\end{equation}
If exact search does not finish, the gap is unknown. A returned legal candidate still supplies the upper bound $\delta\leq c(\widehat\gamma)$.

\subsection{Exact Search, Replay, and Certificates}
\label{sec:exactprelim}

The shortest path view explains why finding an alignment and checking one have different computational demands~\cite{Zelst2018}. A search state can be represented as $(i,m)$, where $i$ events have been consumed and $m$ is the current marking. Synchronous and log moves advance $i$; synchronous and model moves update $m$. Search starts at $(0,m_0)$ and ends at $(n,m_f)$. Equivalently, one constructs a synchronous product of the Petri net of the process and a linear Petri net for the trace. The path cost is the alignment cost.

A* prioritizes states by the sum of the cost already incurred and an estimate of the remaining cost. An estimate is \emph{admissible} when it never overestimates the true minimum remaining cost. Marking equation relaxations obtain useful lower bounds by requiring token balance while relaxing parts of the firing order constraints~\cite{Dongen2018}. They can guide an exact search because their relationship to the original problem is formally justified. LARA's neural scores have no such property and are never substituted for admissible bounds.

Replay solves a simpler task: given a proposed sequence, check each move in order. It confirms the labels, enabledness, complete trace consumption, and final marking, then sums the stated costs. It does not search for an alternative. The amount of work follows the supplied sequence and its incident arcs rather than the full space of alternative executions. This makes replay a practical check even when finding the best sequence is difficult.

There are therefore three distinct levels of evidence. A raw proposal may fail replay. A proposal that passes replay is an executable explanation with a verified upper bound. Finally, if an exact backend proves that its cost equals $\delta$, it is an optimal explanation. Nonnegative costs also imply that any legal zero cost alignment is mathematically optimal; the present certification implementation nevertheless calls the exact backend in certified mode. We keep this implementation fact separate from the mathematical observation.

\subsection{Machine Learning Terminology for Process Mining Readers}
\label{sec:mlprelim}

The purpose of the neural component is to reuse patterns from earlier solved alignment problems. In \emph{supervised machine learning}, a dataset contains inputs and desired outputs. Here the input is a Petri net--trace pair, and the targets come from a replay checked exact alignment. A neural network is a parameterized numerical function $f_\theta$, where $\theta$ collects its adjustable weights. Training changes these weights to reduce a \emph{loss}: a numerical penalty for disagreeing with the targets. After training, \emph{inference} evaluates $f_\theta$ with fixed weights on a new input. A saved set of weights and its architecture settings is a \emph{checkpoint}.

An \emph{embedding} is a vector of learned numerical features. A transition embedding can encode aspects of its label, incoming and outgoing structure, and relationship to the observed trace. Its coordinates need not have individually interpretable meanings. Embeddings let the same calculations operate on Petri nets of different sizes: one vector is created per place, transition, or event, and the parameters are shared across them. The number of objects may change without allocating a separate model for each Petri net.

\emph{Message passing} updates a graph node using information from its neighbors. Typed messages can distinguish a place feeding a transition from a place receiving its output. \emph{Attention} instead lets an object combine information from a collection of other objects using input dependent weights~\cite{Vaswani2017,Gilmer2017}. This is helpful when the next event is ambiguous but a distant suffix event identifies a likely branch. Attention gives access to context; it does not execute the Petri net or enforce its token rules.

For clarity, a single attention operation can be written as
\begin{equation}
 \Att(Q,K,V)=\softmax\!\left(\frac{QK^\top}{\sqrt{d_k}}\right)V.
 \label{eq:attention}
\end{equation}
Rows of $Q$ are \emph{queries}, representing objects seeking information. Rows of $K$ are \emph{keys}, used to score their relevance, and rows of $V$ are the information vectors to combine. The softmax acts row by row: for a vector $z$, $\softmax(z)_j=e^{z_j}/\sum_k e^{z_k}$. Its entries are nonnegative and sum to one, so each output is a weighted combination of value vectors. The scale $\sqrt{d_k}$ controls the magnitude of the dot products. These arrays are learned linear projections of embeddings, rather than externally supplied queries or keys.

In \emph{self attention}, queries, keys, and values originate from the same collection, such as all events. In \emph{cross attention}, they originate from two collections, such as events querying transitions. \emph{Multihead attention} computes several such combinations using different projections and combines their outputs. A \emph{transformer block} adds a small feedforward neural network, residual additions, and layer normalization. Residual additions retain the earlier representation while adding an update; layer normalization rescales a vector's coordinates to stabilize training. A feedforward neural network applies learned affine maps and a nonlinear activation to each vector. None of these operations supplies a formal conformance guarantee.

The final output score before any probability conversion is a \emph{logit}. A larger logit means that the neural network prefers that option under its learned scoring function. Softmax can normalize logits for training, but neither a raw logit nor a normalized score is automatically a calibrated probability that a complete alignment is optimal. A \emph{mask} excludes incompatible options from a calculation. LARA uses exact label equality for one mask and explicit token semantics for enabledness; these are different restrictions.

Training uses three disjoint roles for data. The training split updates $\theta$; the validation split selects settings or a checkpoint; the test split estimates performance after that selection. A full pass through the training split is an \emph{epoch}. A \emph{batch} groups losses before a parameter update. Gradient based optimization computes how changing parameters would affect the loss, then adjusts them. \emph{Regularization} discourages brittle solutions; examples here include weight decay, dropout (temporarily suppressing random features during training), and label smoothing (softening a target that would otherwise be fully certain). An \emph{ablation} disables a component while holding the surrounding procedure fixed, so its contribution can be assessed.

\subsection{Pretraining, Reuse, and Distribution Shift}

Reusable models matter here because a new process should not require another exact labeled training campaign before candidate generation can begin. \emph{Pretraining} fits the scorer on many synthetic Petri net--trace pairs. \emph{Fine tuning} would subsequently update its weights on a target process. LARA's transfer experiments use the pretrained checkpoint without fine tuning, which we call zero shot application with respect to that target process. It still receives the target Petri net and trace as inputs.

Foundation models are commonly associated with broad pretraining and adaptation to a range of downstream tasks~\cite{Bommasani2021}. Process specific proposals argue that event data and process representations warrant their own reusable models~\cite{Rizk2022}. LARA explores the reuse principle within one task. Its approximately 1.7 million parameters, small synthetic corpus, fixed costs, and limited training motifs do not demonstrate broad process mining competence. It is also not a language model: a label such as ``approve application'' is treated as an equality identifier, without tokenizing or interpreting its words.

\emph{Distribution shift} occurs when deployment inputs differ from training inputs. Relevant shifts include longer traces, larger Petri nets, new structural motifs, or different label vocabularies. An exact semantics check continues to have the same meaning under such shifts, but learned rankings may become worse. Evaluating reuse therefore requires separate measurements of candidate legality, optimality, and latency. Acceptance by replay demonstrates successful execution of the candidate, not statistical generalization of the scorer.

\subsection{Research Questions}
\label{sec:rqs}

The definitions above make it possible to ask questions that distinguish the predictor's value from the guarantees of its wrapper. The study addresses three research questions:
\begin{enumerate}[label=\textbf{RQ\arabic*.},leftmargin=*]
\item \textbf{Candidate validity and quality.} How often does candidate generation produce a complete legal alignment, how often is its cost optimal, and how large are its remaining cost gaps? Replay acceptance, exact cost agreement, and gap distributions provide the evidence in Sections~\ref{sec:indist} and~\ref{sec:failure}.
\item \textbf{Contribution of machine learning.} When do learned rankings improve the same decoder compared with deterministic structural choices? A paired ablation and representation specific analysis isolate effects on duplicate labels and identifier sensitivity in Sections~\ref{sec:ablation} and~\ref{sec:representations}.
\item \textbf{Reuse and practical efficiency.} How do quality and running time change across representations, scales, real event logs, and established approximation methods? Sections~\ref{sec:representations}--\ref{sec:real} assess these changes while keeping the pretrained parameters fixed.
\end{enumerate}
These questions do not assume that a neural candidate is preferable to exact search. A useful answer may identify inputs for which its inference cost buys little or no improvement. Nor does a high retained candidate rate answer whether certification is faster: that requires measuring a solver that actually uses the candidate to reduce exact work.

\subsection{Contributions and Scope}
\label{sec:contributions}

The contributions are organized around an implementable alignment pipeline and evidence about its limits:
\begin{enumerate}[label=\textbf{C\arabic*.},leftmargin=*]
\item \textbf{A reusable Petri net--trace scorer.} One neural network for graphs and sequences predicts transition aware move preferences for variable size inputs and previously unseen label vocabularies, without training a neural network per target Petri net (Section~\ref{sec:scorer}).
\item \textbf{An explicit verification contract.} The scorer is integrated with decoding that tracks the marking, independent replay, and exact cost comparison. The resulting output states distinguish a proposal, a verified upper bound, a certified candidate, and an exact replacement (Section~\ref{sec:assurance}). This contract can also wrap candidates produced without a neural network.
\item \textbf{A controlled synthetic supervision procedure.} Behavior families pair different representations of the same visible behavior, reuse corrupted observations across the pair, and obtain concrete transition targets from exact alignments (Section~\ref{sec:training}).
\item \textbf{An empirical attribution and transfer study.} The evaluation compares guided and unguided decoding, representation pairs, larger synthetic inputs, existing approximations, and real event logs. It reports where machine learning helps, where it has no measured effect, and where competing methods perform better (Section~\ref{sec:evaluation}).
\end{enumerate}
The work concerns completed control flow traces on labeled Petri nets with specified initial and final markings. It does not establish guarantees for data aware or object centric alignments, online prefix prediction, arbitrary cost policies, or learned lower bounds. The implementation described in Section~\ref{sec:implementation} makes these contributions inspectable through command line experiments and an interactive workbench.
